%%%%%%%%%%%%%%%%%%%%%%%%%%%%%%%%%%%%%%%%%%%%%%%%%%%%%
%%% Task 1 %%%%%%%%%%%%%%%%%%%%%%%%%%%%%%%%%%%%%%%%%%%%
%%%%%%%%%%%%%%%%%%%%%%%%%%%%%%%%%%%%%%%%%%%%%%%%%%%%%
\task{Transistor-based AC/DC converters}

\taskGerman{Transistorbasierter AC/DC-Umrichter}
A city transport company has installed a \SI{500}{\kilo\watt} charging station for electric buses.
The charger uses a three-phase active front-end (AFE) converter, shown in Fig. \ref{fig:Three-phase_AFE}, to convert power from the
utility grid into a controlled DC supply for the bus battery.
The charging station is connected to a three-phase \SI{400}{\volt} RMS line-to-line, \SI{50}{\hertz} utility
supply. Assume an average DC link voltage of \SI{700}{\volt}, that the efficiency is \SI{96}{\percent}, and that the pulse width modulated switching frequency is set to \SI{10}{\kilo\hertz} .

\begin{germanblock}
	Ein städtisches Verkehrsunternehmen hat eine \SI{500}{\kilo\watt} Ladestation für Elektrobusse installiert.
	Das Ladegerät verwendet einen dreiphasigen AFE-Stromrichter (Active Front-End), um die Energie
	aus dem Versorgungsnetz in eine geregelte Gleichstromversorgung für die Busbatterie umzuwandeln. Die Ladestation ist an ein dreiphasiges Versorgungsnetz mit \SI{400}{\volt}
	(Effektivwert der verketteten Spannung) und \SI{50}{\hertz} angeschlossen. Nehmen Sie an, dass die Spannung am DC-Link bei \SI{700}{\volt} im Mittel liegt, dass der Wirkungsgrad des Stromrichters \SI{96}{\percent} beträgt und dass die Schaltfrequenz der Pulsweitenmodulationauf \SI{10}{\kilo\hertz} eingestellt ist.
\end{germanblock}
\begin{figure}[htb]
	\centering
	\begin{circuitikz}[]
		\draw (0,4) coordinate (A) to [capacitor, *-*, v = $\frac{u_2(t)}{2}$, voltage = straight] ++(0,-3) coordinate (Z) to [capacitor, *-*, v = $\frac{u_2(t)}{2}$, voltage = straight] ++(0,-3) coordinate (B)
			(Z) node[rground, rotate=-90, name = gnd1]{}
			(A) to [short, *-o] ++(-2,0) coordinate (Y)
			(B) to [short, *-o] ++(-2,0) coordinate (X)
			(Y) to [open, o-o, v = $u_2(t)$, voltage shift= 3pt, voltage = straight] (X)
			(A) to [short, i<=$i_\mathrm{conv}(t)$] ++(2,0) coordinate (E)
			to [Tnpn, n=npn1, invert, bodydiode] ++(0,-2) coordinate (C)
			to [short, *-] ++(1,0) to [crossing] ++(2,0) to [crossing] ++(2,0)
			to [short, i<=$i_{1\mathrm{a}}(t)$] ++(1,0) coordinate (U) to [short] ++(1,0) coordinate (G)
			(C) to [short] ++(0,-2)
			to [Tnpn, n=npn2, invert, bodydiode] ++(0,-2) coordinate (D)
			(E) to [short, *-] ++(2,0) coordinate (I)
			to [Tnpn, n=npn3, invert, bodydiode] ++(0,-2)
			to [short] ++(0,-1) coordinate (F)
			to [short] ++(0,-1)
			to [Tnpn, n=npn4, invert, bodydiode] ++(0,-2)
			to [short, -*] ++(-2,0)
			to [short, -o] (B)
			(I) to [short, *-] ++(2,0)
			to [Tnpn, n=npn5, invert, bodydiode] ++(0,-2) coordinate (J)
			to [short] ++(0,-2) coordinate (K)
			to [Tnpn, n=npn6, invert, bodydiode] ++(0,-2)
			to [short, -*] ++(-2,0)
			(F) to [short, *-] ++(1,0) to [crossing] ++(2,0) to [short, i<=$i_{1\mathrm{b}}(t)$] ++(1,0) coordinate (V) to [short] ++(1,0) coordinate (H)
			(K) to [short, *-] ++(1,0) to [short, i<=$i_{1\mathrm{c}}(t)$] ++(1,0) to [short] ++(0.5,0) coordinate(W) to [short] ++(0.5,0) coordinate (L)
			(G) to [open, o-o, v^=$u_{1\mathrm{ab}}(t)$, voltage shift=3pt, voltage = straight] (H)
			(H) to [open, o-o, v^=$u_{1\mathrm{bc}}(t)$, voltage shift=3pt, voltage = straight] (L)
		($(L) + (2,0)$) to [open, v_=$u_{1\mathrm{ca}}(t)$, voltage shift=3pt, voltage = straight] ($(G) + (2,0)$);
		\draw[->] (A) -- ++(-1,0) node[midway, above, xshift=-0.5cm] {$I_2 \approx \mathrm{const.}$};
		\draw[->] (A) -- ++(0,-0.75) node[midway, right] {$i_\mathrm{C}(t)$};
		\draw let \p1 = (npn1.B) in node[anchor=east] at (\x1,\y1) {$T_1$};
		\draw let \p1 = (npn2.B) in node[anchor=east] at (\x1,\y1) {$T_2$};
		\draw let \p1 = (npn3.B) in node[anchor=east] at (\x1,\y1) {$T_3$};
		\draw let \p1 = (npn4.B) in node[anchor=east] at (\x1,\y1) {$T_4$};
		\draw let \p1 = (npn5.B) in node[anchor=east] at (\x1,\y1) {$T_5$};
		\draw let \p1 = (npn6.B) in node[anchor=east] at (\x1,\y1) {$T_6$};
		\draw
		(W |- D) node[rground, anchor = south, name = gnd2]{};
		\draw[->] ($(W) + (0.5,-0.2)$) to ($(W |- gnd2) + (0.5,0.2)$);
		\draw[->] ($(V) + (0.5,-0.2)$) to ($(V |- gnd2) + (0.5,0.2)$);
		\draw[->] ($(U) + (0,-0.2)$) to ($(U |- gnd2) + (0,0.2)$);
		\draw ($(W)!0.5!(gnd2) + (0.5,0)$) node[anchor = west] {$u_{1i0}(t)$};
	\end{circuitikz}
	\caption{Three-phase two-level AC/DC converter}
	\label{fig:Three-phase_AFE}
\end{figure}
\subtask{During commissioning, you observed that the PWM controller currently operates with a modulation index of 0.90.
	Check if the present modulation is sufficient to generate the required AC voltage for grid-connected operation.}{2}
\subtaskGerman{Während der Inbetriebnahme haben Sie festgestellt, dass der PWM-Regler derzeit mit einem Modulationsgrad von 0,90 arbeitet. Prüfen Sie, ob die vorliegende Modulation ausreicht, um die für den Netzparallelbetrieb erforderliche Wechselspannung zu erzeugen.}
\begin{solutionblock}
	For the two-level converter, the fundamental peak phase voltage is given by:
	$$\hat{u}_1^{(1)} = \frac{mU_2}{2}.$$
	The required AC phase-voltage peak is obtained from the line-to-line RMS voltage $U_{1,\mathrm{LL}}$:
	$$\hat{u}_1 = \sqrt{2}\,\frac{U_{1,\mathrm{LL}}}{\sqrt{3}}.$$
	Hence, the modulation index required is:
	$$m = \frac{2\sqrt{2}U_{1,\mathrm{LL}}}{\sqrt{3}U_2}
		\approx \frac{2\sqrt{2}\cdot \SI{400}{\volt}}{\sqrt{3}\cdot \SI{700}{\volt}}
		\approx 0.933 > 0.90.$$
	Therefore, the present modulation is not sufficient for grid-connected operation.
\end{solutionblock}
\subtask{Briefly explain how increasing the switching frequency $f_\mathrm{s}$ affects the grid current harmonic distortion and switching losses.}{1}
\subtaskGerman{Erläutern Sie kurz, wie sich eine Erhöhung der Schaltfrequenz auf die Oberschwingungsverzerrung der Netzströme und die Schaltverluste auswirkt.}
\begin{solutionblock}
	The higher $f_\mathrm{s}$ the better the approximation of a sine wave, the lower the current harmonic distortion, but the higher the switching losses.

\end{solutionblock}
\subtask{Assume $E_\mathrm{on} = \SI{2.8}{\milli\joule}, E_\mathrm{off} = \SI{3.6}{\milli\joule}$ for the loss work per on/off switch instant.
	Calculate the average switching losses at a switching frequency of  $f_\mathrm{s}=\SI{10}{\kilo\hertz}$.}{1}
\subtaskGerman{Unter der Annahme von $E_\mathrm{on} = \SI{2.8}{\milli\joule}, E_\mathrm{off} = \SI{3.6}{\milli\joule}$ pro Schaltvorgang sind die mittleren Schaltverluste bei einer Schaltfrequenz von $f_\mathrm{s}=\SI{10}{\kilo\hertz}$ zu berechnen.}
\begin{solutionblock}
	The power loss per switch is given by
	$$P_\mathrm{sw} = f_\mathrm{s} \cdot (E_\mathrm{on} + E_\mathrm{off}),$$
	and the total switching loss is
	$$P_\mathrm{total} = 6\cdot P_\mathrm{sw} = \SI{384}{\watt}.$$

\end{solutionblock}
\subtask{For nominal operation at \SI{500}{\kilo\watt} battery charging, assess the following:
	\renewcommand{\labelenumi}{\alph{enumi})}
	\begin{enumerate}
		\item Calculate the RMS grid current assuming unity power factor.
		\item What is the effect on the RMS grid current when the power factor reduces to 0.95 (lagging)?
		\item Determine the reactive power exchanged with the grid for a power factor of 0.95 (lagging). State whether the converter is absorbing or supplying reactive power.
	\end{enumerate}}{3}
\subtaskGerman{Für den Nennbetrieb bei \SI{500}{\kilo\watt} Batterieladeleistung sind die folgenden Punkte zu bewerten:
	\renewcommand{\labelenumi}{\alph{enumi})}
	\begin{enumerate}
		\item Berechnen Sie den Effektivwert des Netzstroms  unter der Annahme eines Leistungsfaktors von eins.
		\item Welchen Einfluss hat eine Reduzierung des Leistungsfaktors auf 0,95 (induktiv) auf den Effektivwert des Netzstroms?
		\item Bestimmen Sie die mit dem Netz ausgetauschte Blindleistung für einen Leistungsfaktor von 0,95 (induktiv). Geben Sie an, ob der Stromrichter Blindleistung aufnimmt oder abgibt.
	\end{enumerate}}
\begin{solutionblock}
	\renewcommand{\labelenumi}{\alph{enumi})}
	\begin{enumerate}
		\item The active power drawn from the grid is
		      $$P_{\mathrm{g}} = \frac{P_{\mathrm{dc}}}{\eta} = \frac{\SI{500}{\kilo\watt}}{0.96} = \SI{520.83}{\kilo\watt}.$$
		      For unity power factor, the three-phase power equation is
		      $$P_{\mathrm{g}} = \sqrt{3}U_{1,\mathrm{LL}}I_1\cos\phi,$$
		      with $\cos\phi = 1$. Therefore,
		      $$I_{1,\mathrm{PF}=1} = \frac{P_{\mathrm{g}}}{\sqrt{3}U_{1,\mathrm{LL}}}
			      = \frac{\SI{520.83}{\kilo\watt}}{\sqrt{3}\cdot \SI{400}{\volt}}
			      = \SI{751.8}{\ampere}.$$
		\item For a lagging power factor of $0.95$,
		      $$\phi = \cos^{-1}(0.95) = \SI{18.19}{\degree}.$$
		      The line current becomes
		      $$I_{1,\mathrm{PF}=0.95} = \frac{P_{\mathrm{g}}}{\sqrt{3}U_{1,\mathrm{LL}}\cos\phi}
			      = \SI{791.3}{\ampere}.$$
		      Thus, a lower power factor increases the RMS current drawn from the grid.
		\item The reactive power is
		      $$Q_1 = \sqrt{3}U_{1,\mathrm{LL}}I_{1,\mathrm{PF}=0.95}\sin\phi
			      = \SI{171.4}{\kilo\volt\ampere}.$$
		      Because the power factor is lagging, the converter absorbs reactive power from the grid.
	\end{enumerate}
\end{solutionblock}
\subtask{The charger's DC-link has a total of \SI{90}{\milli\farad} capacitance and its DC voltage ripple needs to be checked. Hence, derive an (approximate) worst-case expression for the peak-to-peak voltage ripple across the DC-link capacitor assuming a constant load.}{4}
\subtaskGerman{Der DC-Link des Ladegeräts verfügt über eine Gesamtkapazität von \SI{90}{\milli\farad} und dessen Spannungsschwankung muss überprüft werden. Leiten Sie daher eine (ungefähre) Worst-Case-Abschätzung für die Spitzenspannungsschwankung über dem DC-Link-Kondensator ab unter der Annahme einer konstanten Last.}
\begin{solutionblock}
	Compared to a single-phase AFE rectifier, the three-phase AFE has no voltage ripple due to the fundamental frequency grid voltage zero-crossings since the fundamental power flow is continuous. Hence, the DC-link voltage ripple is due to the PWM switching only. The DC-side converter current of a three-phase two-level converter can be written as
	\begin{equation*}
		i_{\mathrm{conv}}(t)
		=
		q_\mathrm{a}(t)i_\mathrm{a}(t)+q_\mathrm{b}(t)i_\mathrm{b}(t)+q_\mathrm{c}(t)i_\mathrm{c}(t),
		\qquad q_{\mathrm{a},\mathrm{b},\mathrm{c}}\in\{0,1\}.
	\end{equation*}
	Since the three-phase currents satisfy
	\begin{equation*}
		i_\mathrm{a}(t)+i_\mathrm{b}(t)+i_\mathrm{c}(t)=0,
	\end{equation*}
	the possible DC-side currents are
	\begin{equation*}
		i_{\mathrm{conv}}(t)
		\in
		\{0,\pm i_\mathrm{a}(t),\pm i_\mathrm{b}(t),\pm i_\mathrm{c}(t)\}.
	\end{equation*}
	Hence, a conservative upper bound is
	\begin{equation*}
		\left|i_{\mathrm{conv}}(t)\right|
		\leq \hat i_{\mathrm{1}}.
	\end{equation*}
	The average DC-link load current is
	\begin{equation*}
		I_{\mathrm{2}}
		=
		\frac{P_{\mathrm{dc}}}{U_{\mathrm{2}}}
		=
		\frac{500\,\mathrm{kW}}{700\,\mathrm{V}}
		\approx \SI{714.3}{\ampere}.
	\end{equation*}
	Using the result from Task 1.4 for unity power factor,
	\begin{equation*}
		I_{\mathrm{1}}
		=
		\frac{P_{\mathrm{dc}}/\eta}
		{\sqrt{3}U_{\mathrm{1,LL}}}
		\approx \SI{751.8}{\ampere},
	\end{equation*}
	and therefore
	\begin{equation*}
		\hat i_{\mathrm{1}}
		=
		\sqrt{2}I_{\mathrm{1}}
		\approx \SI{1063.1}{\ampere}.
	\end{equation*}
	The capacitor current is
	\begin{equation*}
		i_\mathrm{C}(t)=i_{\mathrm{conv}}(t)-I_{\mathrm{2}},
	\end{equation*}
	such that the worst-case magnitude can be bounded by
	\begin{equation*}
		|i_\mathrm{C}(t)|
		\leq
		\hat i_{\mathrm{1}}+I_{\mathrm{2}}
		\approx \SI{1777.4}{\ampere}.
	\end{equation*}
	For a worst-case estimate, assume that this current acts over one complete switching period
	\begin{equation*}
		T_\mathrm{s}=\frac{1}{f_\mathrm{s}}.
	\end{equation*}
	With
	\begin{equation*}
		i_\mathrm{C}=C\frac{\mathrm{d}u_{\mathrm{2}}}{\mathrm{d}t},
	\end{equation*}
	the peak-to-peak switching ripple is therefore bounded by
	\begin{equation*}
		\Delta U_{\mathrm{2,pp}}
		\leq
		\frac{|i_C|_{\max}T_\mathrm{s}}{C}
		=
		\frac{\hat i_{\mathrm{1}}+I_{\mathrm{2}}}
		{C f_\mathrm{s}}.
	\end{equation*}
	Insertion of the given values yields
	\begin{equation*}
		\Delta U_{\mathrm{2,pp}}
		\leq
		\frac{\SI{1777.4}{\ampere}}
		{\SI{90}{\milli\farad}\cdot \SI{10}{\kHz}}
		\approx \SI{1.97}{\volt}.
	\end{equation*}
	Thus, the relative peak-to-peak ripple is bounded by
	\begin{equation*}
		\frac{\Delta U_{\mathrm{2,pp}}}{U_{\mathrm{2}}}\,100\%
		\leq
		\frac{\SI{1.97}{\volt}}{\SI{700}{\volt}}\,100\%
		\approx 0.28\%.
	\end{equation*}
\end{solutionblock}
